\begin{figure}[h]
\centering
\begin{tikzpicture}[scale=0.8]
% 左图:模糊prompt
\begin{scope}[xshift=-3.5cm]
\node at (0, 1.5) {\textbf{\ifdefined\ENGLISH Fuzzy Hook\else 模糊的钩子\fi}};
\draw[thick, ->] (0, 1.2) -- (0, 0.5);
\node[draw, circle, minimum size=2cm, fill=paleGray] (neuron1) at (0, -0.8) {};
\node[font=\tiny] at (-0.5, -0.5) {\ifdefined\ENGLISH cat\else 猫\fi};
\node[font=\tiny] at (0.5, -0.5) {\ifdefined\ENGLISH dog\else 狗\fi};
\node[font=\tiny] at (-0.5, -1.1) {\ifdefined\ENGLISH fluffy\else 毛茸茸\fi};
\node[font=\tiny] at (0.5, -1.1) {\ifdefined\ENGLISH pet\else 宠物\fi};
\draw[gray, thick, dashed] (0, -0.8) circle (0.9cm);
\node at (0, -2.5) {\small \ifdefined\ENGLISH Activates multiple concepts\else 激活多个概念\fi};
\node at (0, -3) {\small \ifdefined\ENGLISH $\to$ Mixed output\else → 输出混杂\fi};
\end{scope}

% 右图:精准prompt
\begin{scope}[xshift=3.5cm]
\node at (0, 1.5) {\textbf{\ifdefined\ENGLISH Precise Hook\else 精准的钩子\fi}};
\draw[thick, ->] (0, 1.2) -- (-0, 0.5);
\node[draw, circle, minimum size=2cm, fill=paleGray] (neuron2) at (0, -0.8) {};
\node[font=\tiny] at (-0.5, -0.5) {\ifdefined\ENGLISH cat\else 猫\fi};
\node[font=\tiny] at (0.5, -0.5) {\ifdefined\ENGLISH dog\else 狗\fi};
\node[font=\tiny] at (-0.5, -1.1) {\ifdefined\ENGLISH fluffy\else 毛茸茸\fi};
\node[font=\tiny] at (0.5, -1.1) {\ifdefined\ENGLISH pet\else 宠物\fi};
\draw[darkCyan, thick] (-0.5, -0.5) circle (0.25cm);
\node at (0, -2.5) {\small \ifdefined\ENGLISH Activates target precisely\else 精确激活目标\fi};
\node at (0, -3) {\small \ifdefined\ENGLISH $\to$ Clear output\else → 输出清晰\fi};
\end{scope}
\end{tikzpicture}
\caption{\ifdefined\ENGLISH Relationship between hook precision and concept activation\else 钩子精度与概念激活的关系\fi}
\label{fig:hook-precision}
\end{figure}
